\documentclass[12pt]{article}
\usepackage{setspace}
\usepackage{amsmath,amssymb,euscript,graphicx}
\usepackage{xurl}
\usepackage{graphicx}
\usepackage{subfigure}
\usepackage{wrapfig}
\usepackage[compact]{titlesec}
\usepackage{fancybox,color}
\usepackage[footnotesize]{caption}
\usepackage{cite}
\usepackage{enumitem}
\usepackage[normalem]{ulem}
\usepackage[top=0.9375in, right=0.9375in, left=0.9375in, bottom=0.9375in, centering]{geometry}
\usepackage{bm}

\usepackage{amsthm}


% -- Loading the code block package:
\usepackage{listings}
% -- Basic formatting
\usepackage[utf8]{inputenc}
\usepackage[english]{babel}
\usepackage{times}
\setlength{\parindent}{8pt}
\usepackage{indentfirst}
% -- Defining colors:
\usepackage[dvipsnames]{xcolor}
\definecolor{codegreen}{rgb}{0,0.6,0}
\definecolor{codegray}{rgb}{0.5,0.5,0.5}
\definecolor{codepurple}{rgb}{0.58,0,0.82}
\definecolor{backcolour}{rgb}{0.95,0.95,0.92}
% Definig a custom style:
\lstdefinestyle{mystyle}{
    backgroundcolor=\color{backcolour},   
    commentstyle=\color{codepurple},
    keywordstyle=\color{NavyBlue},
    numberstyle=\tiny\color{codegray},
    stringstyle=\color{codepurple},
    basicstyle=\ttfamily\footnotesize\bfseries,
    breakatwhitespace=false,         
    breaklines=true,                 
    captionpos=t,                    
    keepspaces=true,                 
    numbers=left,                    
    numbersep=5pt,                  
    showspaces=false,                
    showstringspaces=false,
    showtabs=false,                  
    tabsize=2
}
% -- Setting up the custom style:
\lstset{style=mystyle}


% IEEE uses Times Roman font, so we'll default to Times.
% These three commands make up the entire times.sty package.
\renewcommand{\sfdefault}{phv}
\renewcommand{\rmdefault}{ptm}
\renewcommand{\ttdefault}{pcr}

\newcommand{\squeeze}{\vspace{-0.1in}}
\newcommand{\changed}[1]{\textcolor{blue}{#1}} %generates text in blue

%%%%%%%%%%%%%%%%%%%%%%%%%%%%%%%%%%%%%%%%%%%%%%%%%%%%%%%%%%%%%%%%%%%%%%%

\begin{document}

\begin{center}
  {\Large {\bf Lab 0: Setting up your workstation}}
  
\end{center}

\vspace{5pt}
\hrule
\vspace{5pt}

\section{Summary}
For our first lab we will set up your computer to run the necessary software. 
This requires installing Python and installing Microsoft Visual Studio Code.
I have included instructions for both the Mac OS and Windows. 
At the end of this lab you should be able to open a text editor and run a basic python program in the Terminal.

The assignment instructions come directly from the Introduction to Computer Science Course instructions from the Colby College Computer Science Department (CS 151 and CS 152).
Direct credit and acknowledgements of sources are at the end of this document. 

\section{Installing Python}
Python is an interpreted programming language that is commonly used in introduction to computer science classes due to the highly readable syntax.
It is free to install and runs on many different operating systems. 

\subsection{macOS}
Installation Instructions: 
\begin{enumerate}
\item Download Python 3.13.7 at \url{https://www.python.org/downloads/macos/}
\item Open the downloaded file called \textit{python-3.13.7-macos11.pkg} which should be in the Downloads folder. Note if you do not see the \textit{.pkg} this is fine. 
\item Double click on the file and a pop-up window will appear, follow the steps clicking Continue button 3 times and Agree when prompted
\item Click continue one more time
\item Click install when prompted
\item You will enter your computer password and Python will install (this might take a minute or so be patient). After this process finishes, close the installer and delete the file in the Downloads folder.
\end{enumerate}

Check that it is installed correctly:
\begin{enumerate}
\item Open a Terminal, this can be found in the Applications/Utilities folder in Finder or searching for it using Spotlight (Cmd+spacebar, type Terminal and then press Enter)
\item A small black window will open. Type: \texttt{python3 --version} and hit Enter. It should return Python 3.13.7, which means things are installed correctly!
\end{enumerate}

We now want to make it so that when you type the command python it works instead of needing to type python3.
In the Terminal, do the following:
\begin{enumerate}
\item echo \$SHELL
\item If the output is bash:
  \begin{enumerate}
  \item \texttt{sudo nano \textasciitilde{}/.bash\textunderscore profile}
  \item Type your computer password into the terminal
  \item Add the following line: \texttt{alias python=python3}
  \item Add the following line: \texttt{alias pip=pip3}
  \item Save file by hitting Ctrl + x and hitting Y and then Enter (or return) when prompted to save
  \item Close out of the Terminal
  \end{enumerate}
\item If the output is zsh:
  \begin{enumerate}
  \item \texttt{sudo nano \textasciitilde{}/.zshrc}
    \item Type your password when prompted
  \item Add the following line: \texttt{alias python=python3}
  \item Add the following line: \texttt{alias pip=pip3}
  \item Save file by hitting Ctrl + x and hitting Y and then Enter (or return) when prompted to save
  \item A new bar will appear at the bottom with TAB hit enter and everything should save.
  \item Close out of the Terminal
  \end{enumerate}
\item Test that this is working by opening a new Terminal and typing \texttt{python --version}
\end{enumerate}


\subsection{Windows}
Installation Instructions: 
\begin{enumerate}
\item Download Windows installer(64-bit) at

  \url{https://www.python.org/downloads/windows/}

\item Run the installer file called \textit{python-3.13.7-amd64.exe}
\item An installation window will pop up, check the last box at the bottom of the window: ``Add Python 3.13.7 to PATH'' and ``Use admin privileges when installing py.exe''
\item Click ``Install Now'' at the top.
  \item Hit close button once the successful install button appears.
\end{enumerate}

Check that it is installed correctly: 
\begin{enumerate}
\item Press WinKey + r to open the Run box.
\item Type cmd into the box and press Enter
  \item Type \texttt{python --version}, then press Enter. It should print Python 3.13.7 which means things are installed correctly
\end{enumerate}


\section{Instling Microsoft Visual Studio Code}
For all projects we are going to write code in a text editor. While there are a wide range of text editors out there, I recommend using Microscoft Visual Studio Code (VS Code for short).
This editor has built in features we will use for all the assignments.

\subsection{macOS}
Installation Instructions: 
\begin{enumerate}
\item Download Microscoft VS Code at \url{https://code.visualstudio.com/Download}. Select the Mac OS option.
\item You should see a file with a blue icon called Visual Studio Code in your Downloads folder. Drag it into your Applications folder
\item Open VS Code and create a new file called \textbf{test.py}. To make a new file go to the top bar of the application click File $\rightarrow$ New File, a pop up window will appear and you will be prompted to name the new file.
  \item Once you have made your new file, you should see a pop-up in bottom right corner asking you about installing a Python extension. Click Install
\item VS Code should install it, then you should see Python at the bottom right corner of the screen
\item If a drop-down menu should appear with some options at the top. Then Click the entry that says 3.13.7 followed by /usr/local/bin/python3
\item Dismiss the pop-up window that continually appears about ``Linter pylint''. Click Do not show again.
\end{enumerate}

\subsection{Windows}
Installation Instructions: 
\begin{enumerate}
\item Download Microscoft VS Code at \url{https://code.visualstudio.com/Download}. Select the Windows option
\item Dobule-click the installer file to install it on your computer
\item To integrate Python with VS Code, open and create a new file called \textbf{test.py}. To make a new file go to the top bar of the application click File $\rightarrow$ New File, a pop up window will appear and you will be prompted to name the new file.
\item A pop-up in VS code about a Python extension being recommended will appear. Click Install
\item You should see ``Python 3.13.7'' in the bottom right corner of the editor which means you are ready to go
\end{enumerate}

We will now install an additional package necessary on Windows to do the assignments this semester.
\begin{enumerate}
\item Download Git from \url{https://git-scm.com/download/win}
\item Double click and install Git, this will ask you a series of questions agree to the default settings and install.
\item Open Visual Studio Code and press (Ctrl + \textasciigrave) where  \textasciigrave is the same key as \textasciitilde{}.
\item Open the command palette using (Ctrl + Shift + p)
\item Type - Select Default Profile
\item Select Git Bash from the options
\item Close Visual Studio
\item Reopen Visual Studio and you should see bash in the right bottom corner next to the terminal
\end{enumerate}
These instructions came from:
\url{https://stackoverflow.com/questions/42606837/how-do-i-use-bash-on-windows-from-the-visual-studio-code-integrated-terminal}

\section{Show file extensions}
To ensure you can see which files are code files, we need to make sure the proper settings are setup.

\subsection{macOS}
\begin{enumerate}
\item Go to the Finder preferences. (This is found on the top left corner of the screen and click preferences)
\item Click Advanced
\item Make sure \textit{Show all filename extensions} is checked.
\end{enumerate}


\subsection{Windows}
\begin{enumerate}
\item Open up File Explorer 
\item Go to View
 \item Click \textit{File name extensions} 
\end{enumerate}

In both cases you should see \textit{test.py} with the .py showing in your Finder or Explorer window. 

\section{Checking your setup}
In VS Code, open a terminal. The shorthand is Ctrl + \textasciigrave where  \textasciigrave is the same key as \textasciitilde{}.

If you do not want to use the short hand then: 
\begin{itemize}
\item \textbf{Mac} Go to the Terminal menu on top bar and select New Terminal
\item \textbf{Windows} Go to the 3 dots (...) on the top bar, select Terminal, New Terminal
\end{itemize}
In the next section we will go over the details of what a terminal is and why we should use it etc. 
In the terminal, on both operating systems do the following:

\vspace{5mm}

\textit{python $--$version}

\textit{pip $--$version}

\textit{pip install opencv-contrib-python}

\vspace{5mm}

\subsection{Writing First Python Program}
In the terminal type:

\vspace{5mm}
  \textit{mkdir ES1098}
  
  \textit{cd ES1098}
  
  \textit{code myFirstProgram.py}
  \vspace{5mm}
  
This will make a new directory called ES1098, move you into that directory and open a file called \texttt{myFirstProgram.py}. 

In \textbf{myFirstProgram.py}, type:
\begin{lstlisting}[language=Python]
print(``Hello World'')
\end{lstlisting}

Remember to save your file, this is done by type \texttt{Ctrl + s} or File $\rightarrow$ save.

In the terminal type:
\vspace{5mm}

\textit{python myFirstProgram.py}

\vspace{5mm}
What is the output?

\vspace{5mm}
Congratulations! You have completed lab 0 and the setup for the rest of the semester! We will rely on this backbone for all labs and projets.


\section{Acknowledgments}
Section 2, 3, and 4 of these instructions are from Professor Oliver Layton and Professor Hannen Wolfe in the Colby College CS Department. Section 5 is a combination of the prior sources and including work from Professor Bruce Maxwell and Professor Stephanie Taylor in the Colby College CS Department. 
Sources:
\begin{itemize}
\item \footnotesize{\url{https://cs.colby.edu/courses/S25/cs151/projects/00/Lab00Software.html}}
\item \url{https://cs.colby.edu/courses/S25/cs151/projects/01/Lab01First-Python.html}
\item \url{https://cs.colby.edu/courses/S20/cs152-labs/labs/lab01/}
\item \url{https://stackoverflow.com/questions/42606837/how-do-i-use-bash-on-windows-from-the-visual-studio-code-integrated-terminal}
\end{itemize}


\end{document}
